\begin{abstract}
Task-vector subspace overlap is often interpreted as shared learned structure because it exceeds
an isotropic baseline, although linear-layer updates depend on the activations they read. We
introduce a conditional randomisation that preserves each update's spectrum and blockwise
occupancy in frozen-base activation geometry while randomising cross-task directional agreement;
its expected overlap has a closed form. We pair this null with a positive control that plants
shared directions and measures recovery of the delivered signal, exposing consequential
under-randomisation in our initial implementation. In complete ViT-B measurements, the corrected
activation null reproduces about a quarter of above-isotropic overlap, moving over $20$--$26\%$
across our calibration and blocking grid at the reported subspace size, and the residual stays
positive throughout. Module role moves this far
more than architecture does. For shared-start LoRA a regime-matched factor null reproduces about
$80\%$, but a one-line count from the shared rank already gives $82\%$, so there the admissible
row space, not a learned basis, accounts for nearly all observed overlap. These are attribution
statements conditional on what each null holds fixed, and they leave pair rankings and downstream
merge behaviour unchanged.
\end{abstract}
